
\subsubsection*{4.3.2}
\phantomsection
\label{sga3:VIIB:4.3.2}

``There is no Campbell--Hausdorff formula in characteristic
\(p>0\).'' Let us explain. The functorial isomorphism of \SGARefLink{sga3:VIIA:4.3.1}{4.3.1}
depends on the choice of the bases
\(({}^C x_i)_{i\in I}\). Unlike the situation in \SGARefLink{sga3:VIIB:3.2:prop}{3.2}, in
characteristic \(0\), there is in general no bijection
\[
  \mathbf L(C)\otimes_C\sqrt[p]{0_C}
  \longrightarrow
  \mathscr G_p(\mathbf L)(C)
\]
that is ``functorial both in \(C\) and in \(\mathbf L\).''
The following argument shows, for example, that no such
isomorphism exists when \(k\) is a field containing all
\((p^2-1)\)-st roots of unity.

Choose \(\mathbf L\) so that, for every
\(C\in\operatorname{Alf}/k\), the algebra \(\mathbf L(C)\) is the
abelian \(p\)-Lie algebra over \(C\) generated by an element \(x\)
subject to
\[
  x^{(p^2)}=0.
\]
Thus
\[
  \mathbf L(C)
  =
  Cx\oplus Cx^{(p)},
  \qquad
  U_p\bigl(\mathbf L(C)\bigr)
  \simeq
  C[x]/(x^{p^2}).
\]
\phantomsection
\label{sga3:VIIB:4.3.2:source-notation-note}
\noindent\textit{Source note.}\enspace
The French prints \(k[x]/(x^{p^2})\) in the last display, although
the enveloping algebra here is formed over \(C\). A little later it
also identifies \(\mathscr G_p(\mathbf L)(C)\), whose elements are
group-like, with \(\operatorname{Prim}U_p(\mathbf L(C))\).
The target uses the coefficient ring \(C\) and retains the
group-like set defined in \SGARefLink{sga3:VIIA:4.2.1}{4.2.1}; primitive elements instead
describe its Lie algebra in 4.2.2.

We shall show that the only functorial morphism
\[
  \chi_C:
  \mathbf L(C)\otimes_C\sqrt[p]{0_C}
  \longrightarrow
  U_p\bigl(\mathbf L(C)\bigr)
\]
which is compatible with the endomorphisms of \(\mathbf L\) and
sends \(\mathbf L(C)\otimes_C\sqrt[p]{0_C}\) into
\(\mathscr G_p(\mathbf L)(C)\) is the constant morphism with value
\(1\).

Let
\[
  \psi_C:
  \mathbf L(C)\otimes_C\sqrt[p]{0_C}
  \xrightarrow{\ \sim\ }
  \mathscr G_p(\mathbf L)(C)
\]
be the bijection of \SGARefLink{sga3:VIIA:4.3.1}{4.3.1}, namely
\[
  x\otimes a+x^{(p)}\otimes b
  \longmapsto
  \sum_{0\leq i,j<p}
  a^i b^j x^{\,i+pj}.
\]
Composing \(\chi_C\) with \(\psi_C^{-1}\) gives a morphism,
functorial in \(C\),
\[
\begin{aligned}
  \theta_C:
  \mathbf L(C)\otimes_C\sqrt[p]{0_C}
  &\longrightarrow
  \mathbf L(C)\otimes_C\sqrt[p]{0_C},
  \\
  x\otimes a+x^{(p)}\otimes b
  &\longmapsto
  x\otimes P(a,b)+x^{(p)}\otimes Q(a,b).
\end{aligned}
\]
Functoriality in \(C\) implies that \(P(a,b)\) and \(Q(a,b)\) are
the values at \((a,b)\) of two polynomials
\[
  P,Q\in k[X,Y]
\]
which may be assumed to have degree less than \(p\) in each of
\(X\) and \(Y\).\footnote{\label{sga3:VIIB:nde:138}The editors
expanded the original in the argument that follows.}

Let \(\alpha\in k\), and let \(\ell_\alpha\) be the endomorphism of
the \(p\)-Lie algebra \(\mathbf L\) sending
\[
  x\longmapsto\alpha x,
  \qquad
  x^{(p)}\longmapsto\alpha^p x^{(p)}.
\]
Then \(U_p(\ell_\alpha)\) sends \(x\) to \(\alpha x\), and hence
sends \(x^i\) to \(\alpha^i x^i\) for \(i<p^2\). The following
square commutes:
\begin{center}
  \phantomsection
  \label{sga3:VIIB:4.3.2:naturality-endomorphism-square}
  \begin{tikzcd}[row sep=3.2em,column sep=4.8em]
\mathbf L(C)\otimes_C\sqrt[p]{0_C}
  \arrow[r,"\psi_C"]
  \arrow[d,"{\ell_\alpha(C)\otimes_C\operatorname{id}}"]
  & U_p(\mathbf L(C))
      \arrow[d,"{U_p(\ell_\alpha)(C)}"] \\
\mathbf L(C)\otimes_C\sqrt[p]{0_C}
  \arrow[r,"\psi_C"]
  & U_p(\mathbf L(C))
\end{tikzcd}

\end{center}
The hypothesis
\[
  \chi_C\circ(\ell_\alpha(C)\otimes\operatorname{id})
  =
  U_p(\ell_\alpha)(C)\circ\chi_C
\]
therefore gives
\[
  P(\alpha a,\alpha^p b)=\alpha P(a,b),
  \qquad
  Q(\alpha a,\alpha^p b)=\alpha^p Q(a,b).
\]

Write
\[
  P=\sum_{i,j<p}\lambda_{ij}X^iY^j,
  \qquad
  Q=\sum_{i,j<p}\mu_{ij}X^iY^j,
\]
and take
\[
  C=k[X,Y]/(X^p,Y^p).
\]
If \(\lambda_{ij}\ne0\), then
\[
  \alpha^{\,i-1+pj}=1;
\]
if \(\mu_{ij}\ne0\), then
\[
  \alpha^{\,i+p(j-1)}=1.
\]
Taking \(\alpha\) to be a primitive root of unity of order
\(p^2-1\), we conclude that
\[
  P=\lambda X,
  \qquad
  Q=\mu Y
\]
for some \(\lambda,\mu\in k\). Consequently,
\[
  \chi_C(x\otimes a+x^{(p)}\otimes b)
  =
  \sum_{0\leq i,j<p}
  (\lambda a)^i(\mu b)^j x^{\,i+pj}.
\]

Finally, consider the endomorphism \(f\) of \(\mathbf L\) sending
\[
  x\longmapsto x+x^{(p)}.
\]
Take \(C=k[a]/(a^3)\) and compare the coefficients of \(x^p\) and
\(x^{p+1}\) in
\[
  (\chi_C\circ f(C))(x\otimes a)
  \qquad\text{and}\qquad
  \bigl(U_p(f)(C)\circ\chi_C\bigr)(x\otimes a).
\]
This gives
\[
  \lambda=\mu,
  \qquad
  \lambda^2a^2=0,
\]
and hence \(\lambda=\mu=0\), as claimed.
